\section{Bush --- International Politics and the Spread of Quotas for Women in Legislatures}

\subsection{The puzzle and the central argument}

Sarah Sunn Bush examines the rapid global diffusion of gender quotas for legislative representation. Before the 1970s, only a handful of countries had adopted quotas for women, whereas by the late 2000s more than one hundred countries had some form of quota. The striking feature of this diffusion is that many adopters are developing countries in which women's broader social, economic, and political status remains relatively low. This pattern is difficult to explain with a simple modernization account in which economic development, democratization, and improvements in women's socioeconomic position gradually generate demand for greater political equality. Bush argues instead that international politics has played a central role in spreading gender quotas, especially in developing countries after the Cold War.

The key mechanism is the incorporation of gender equality into international democracy promotion. Since the end of the Cold War, democracy assistance has become a large professional field involving governments, international organizations, foundations, nongovernmental organizations, contractors, and experts. Within this field, gender equality and women's political participation increasingly came to be understood as components of legitimate democratization. Gender quotas therefore acquired international legitimacy as an institutional device appropriate for transitional and developing democracies. Once this happened, domestic leaders faced incentives to adopt quotas even when local social conditions did not strongly favor them.

Bush identifies two principal causal pathways. The first is direct international influence in postconflict settings. Where liberalizing United Nations peace operations and other external actors are deeply involved in writing constitutions and electoral laws, international experts can directly press for gender quotas and support domestic actors who favor them. The second pathway is indirect inducement and signaling. Governments that depend on foreign aid, international investment, diplomatic support, or external legitimacy can adopt quotas to signal that they are committed to democracy, inclusion, and good governance. The signal may be sincere or strategic; the important point is that the international legitimacy of quotas makes their adoption politically useful.

\subsection{What counts as a gender quota?}

Bush distinguishes three major forms of gender quotas. Reserved seats set aside a fixed number or percentage of parliamentary seats for women. Legislative quotas require political parties to nominate a specified share of female candidates. Voluntary party quotas are adopted internally by political parties rather than mandated by the state. The first two are legal quotas because they generally require electoral-law or constitutional provisions. They vary widely in size and enforcement, but a quota of around 30 percent is especially common.

The global pattern is important for the argument. Developed Western democracies have often relied on voluntary party quotas, while developing regions have made much greater use of legal quotas, including reserved seats. Reserved-seat systems are particularly common in Africa, Asia, and the Middle East and North Africa, precisely the regions that tend to score poorly on other indicators of women's status. Thus the international spread of quotas does not resemble simple convergence on one universal institutional model. Instead, different types of quotas have spread unevenly across countries with very different social and political conditions.

Bush therefore focuses on national legal quotas in developing countries. These are the policies most clearly connected to the activities of international democracy promoters and most useful as signals of a government's commitment to inclusion. The relevant empirical question is not why every state in the world adopts the same policy, but why developing and transitional states have become especially likely to adopt legally mandated quotas.

\subsection{Alternative explanations: modernization and world polity theory}

The article first considers domestic structural explanations. Revised modernization theory expects economic development, democratization, women's education and labor-force participation, and changing gender attitudes to produce greater gender equality. From this perspective, richer and more democratic countries, countries with higher levels of female socioeconomic development, and countries with less patriarchal cultural traditions should be more likely to adopt quotas. Muslim-majority countries, given arguments about more patriarchal social structures, should be less likely to do so.

The observed pattern already creates problems for this view. Many quota adopters are poor or only partially democratic, and some are countries where women have relatively weak socioeconomic status. Bush does not deny that development can matter, but she argues that modernization variables do not explain the main post-Cold War diffusion pattern among developing countries.

A second alternative is world polity theory. According to this approach, international organizations and transnational networks disseminate standardized scripts of appropriate modern statehood. Countries linked more closely to global civil society should therefore be more likely to adopt internationally legitimate policies. Applied to gender quotas, this perspective suggests that countries with more international nongovernmental organizations, especially women's organizations, with longer histories of women's suffrage, with ratification of the Convention on the Elimination of All Forms of Discrimination Against Women, and with neighboring countries that have already adopted quotas should be more likely to follow suit.

Bush's argument shares world polity theory's interest in legitimacy and appropriate institutions, but differs in the mechanism and in its predictions. She emphasizes international inequalities and incentives rather than general socialization into a universal script. The influence of quotas is strongest where international actors possess unusual leverage: postconflict countries under international supervision and states materially or reputationally dependent on Western donors and democracy promoters. The result is not universal convergence but a distinctive concentration of legal quota adoption in developing states.

\subsection{The democracy establishment and the international legitimacy of quotas}

Bush uses the term ``democracy establishment'' to describe the network of actors involved in democracy assistance and promotion. It includes government agencies such as USAID and the U.S. State Department; quasi-governmental organizations and foundations; NGOs such as Freedom House and the Carter Center; private contractors; European party foundations; and international organizations including the UN, OSCE, European Union, Council of Europe, and International IDEA. These actors share a professional commitment to democracy promotion and increasingly converge on a common set of recommended institutions and practices.

The democracy establishment combines features of transgovernmental networks, transnational advocacy networks, and epistemic communities. Its members exchange expertise, provide technical assistance, promote professional standards, and share normative beliefs about good democratic governance. A common democracy-assistance package has emerged around elections, rule of law, civil society, governance, local political development, and women's participation.

Several processes helped make gender quotas part of this professional repertoire. First, democracy promoters often possessed principled commitments to gender equality. Second, activists, including transnational women's groups, pressured international organizations to pay greater attention to women's rights. Third, democracy promoters learned from prior experiences in countries where quotas were adopted and transferred those lessons elsewhere. Personnel frequently moved among countries and organizations, facilitating diffusion of policy models. Fourth, aid organizations had bureaucratic incentives to adopt measurable indicators of progress, and the number or percentage of women participating in political institutions was an easily quantifiable output. These processes of professionalization, learning, turnover, and standardization reinforced the belief that gender quotas were an appropriate institution for new democracies.

The broader normative environment also changed. United Nations Security Council Resolution 1325 in 2000 explicitly affirmed women's participation in peace and security. In postconflict reconstruction, gender inclusion increasingly came to be seen as an expected feature of legitimate political institutions. This did not mean that every democracy promoter agreed on the precise design of quotas, but some form of gender mechanism became a widely accepted expectation.

\subsection{Three hypotheses about international incentives}

Bush derives three hypotheses. First, countries hosting a democracy-promoting UN peace operation should be more likely to adopt gender quotas. In postconflict settings, international actors are deeply involved in constitutional design and electoral reform. They provide technical advice, funding, and organizational support, and they frequently empower domestic women's organizations. At the strongest end, this influence can approach international imposition; at a minimum, the dependence of postconflict governments on external assistance gives them strong incentives to follow internationally legitimate models.

Second, higher foreign-aid dependence should increase the likelihood of quota adoption. Leaders who rely heavily on external donors have incentives to demonstrate liberal credentials and good governance. Gender quotas are attractive because they are visible, formal, and relatively easy to present as evidence of inclusion. Adoption can therefore help a government cultivate goodwill and legitimacy even if its broader democratic record is mixed.

Third, countries that invite international election monitors should be more likely to adopt quotas. Bush treats election observation mainly as a proxy for otherwise difficult-to-measure incentives to perform democracy for international audiences. A government willing to invite observers is likely to be sensitive to international reputation, foreign aid, investment, debt relief, and legitimacy. The same motivations should make quotas attractive as a second signal of liberalism.

The cases of Kosovo and Jordan illustrate the mechanisms. In Kosovo, the international administration established a 30 percent quota despite domestic resistance, demonstrating direct influence. In Jordan, King Abdullah II introduced reserved seats for women while cultivating closer aid and political ties with the United States and Europe, illustrating how a quota could serve as a signal of modernization and reform to external audiences even in the presence of domestic opposition.

\subsection{Research design and statistical findings}

The quantitative analysis uses country-year data from 1970 to 2006 and excludes long-term consolidated, advanced industrial democracies because Bush argues that quota adoption there follows different causal processes. The dependent variable records the first adoption of a national legal quota. The main international variables measure the presence of a liberalizing UN peace operation, foreign aid received, and whether the most recent election was internationally observed. Alternative explanations include democracy, Islamic heritage, women's international NGOs, years since female suffrage, and the regional diffusion of quotas.

Bush uses a Cox proportional-hazards event-history model, which is appropriate for studying the timing of relatively rare and usually irreversible policy adoption. The results are much more supportive of the international-incentives argument than of modernization theory. Democracy and Islamic heritage are not significant predictors in the main models. World polity variables perform somewhat better: the number of women's INGOs is positively associated with adoption in a simpler model, but this effect weakens in fuller specifications. Years since women's suffrage and the share of regional quota adopters do not provide strong independent explanations.

The variables linked to international incentives perform substantially better. A liberalizing UN peace operation, an internationally observed election, and foreign-aid dependence are all positively associated with quota adoption in the relevant models. In the full model, internationally observed elections and aid dependence retain important effects, while the peace-operation effect remains substantively large though less precisely estimated. Simulations illustrate the magnitude of these relationships: moving from no liberalizing UN peace operation to having one raises the estimated hazard of quota adoption by roughly 157 percent; moving from an unobserved to an internationally observed election raises it by about 147 percent; and moving foreign aid from the twenty-fifth to the seventy-fifth percentile raises the hazard by roughly 214 percent. These results support the view that external leverage and international signaling matter more than the domestic modernization indicators that conventional theories emphasize.

\subsection{Afghanistan's 2004 quota: process tracing}

Bush complements the statistical analysis with a detailed study of Afghanistan, which adopted a quota in its January 2004 constitution. The Bonn Agreement of 2001 called for a broad-based, gender-sensitive, multiethnic government, and the UN Assistance Mission in Afghanistan became deeply involved in supporting the constitutional process. The resulting constitution reserved 27 percent of seats in the lower house and 17 percent in the upper house for women.

The case study examines three observable implications of the international-incentives theory: whether international experts influenced the constitutional process, whether they regarded quotas as appropriate, and whether international actors empowered Afghan women who advocated them. Evidence supports all three. Although international officials deliberately maintained a ``light footprint'' to avoid the appearance of foreign imposition, expert advisers played an important role in ensuring that the constitution met international standards, particularly regarding women's rights and gender equality. UN officials and other democracy promoters broadly agreed that a postconflict constitution needed some mechanism guaranteeing women's political representation.

International influence was not simply top-down. Afghan women also participated actively. Women made up a significant share of the constitutional commissions and Loya Jirga, and some successfully campaigned to expand the lower-house quota. Yet their capacity to do so was partly produced by international support. UNDP, UNIFEM, U.S. officials, the National Endowment for Democracy, and other organizations funded women's institutions, provided technical expertise, organized meetings with international experts, and helped select women for participation in the constitutional process. In this sense, domestic advocacy and international influence were mutually reinforcing, but the latter supplied crucial resources and leverage.

The case also shows the limits and political ambiguities of externally promoted inclusion. Many Afghan women did not initially demand quotas, and some elite women objected to the idea of receiving political ``charity.'' Later controversies, such as the Shia Family Law, demonstrated that international actors sometimes prioritized symbolic legal reforms over women's security and socioeconomic concerns. Nevertheless, the quota had major consequences: in the 2005 lower-house election, most female members owed their seats to the reserved-seat mechanism. Bush concludes that domestic support supplied local legitimacy, but international pressure and inducements were necessary to tip the political balance in favor of quota adoption.

\subsection{Conclusion and implications}

Bush concludes that the spread of legal gender quotas in developing countries is primarily an international political phenomenon. Quotas are not simply the outcome of modernization or generalized world-polity convergence. They spread because the international democracy-promotion community made them symbols of legitimate democratization and because international actors possessed leverage over particular countries. In postconflict states, this influence could be direct; in aid-dependent and reputation-sensitive states, it worked indirectly through inducement and signaling.

The argument matters substantively because quotas affect representation and, in turn, the policies and attitudes associated with women's political participation. It also matters theoretically because it shows how international actors reshape domestic institutions through a mixture of normative legitimacy and material asymmetry. Norms do not affect all states equally. Internationally legitimate policies may diffuse most rapidly where external actors have the strongest authority, resources, and inducements. The article therefore links constructivist concerns about legitimacy to a more strategic account of international influence, showing how ideas become politically consequential through concrete distributions of power and dependence.

\section{Finnemore and Sikkink --- International Norm Dynamics and Political Change}

\subsection{Purpose and contribution}

Martha Finnemore and Kathryn Sikkink provide one of the foundational constructivist accounts of how international norms emerge, spread, and become internalized. Their starting point is that normative and ideational concerns are not new to international relations. Earlier realists, integration theorists, and scholars of decolonization all treated legitimacy, social purpose, ideology, and identity as politically important. What changed during the behavioral and rational-choice turns of the 1960s through the 1980s was less the disappearance of normative phenomena than the methodological standards applied to their study. The renewed ``ideational turn'' therefore represents a return to old substantive concerns equipped with more explicit causal arguments and more systematic empirical research.

The authors seek to answer several questions: how can researchers identify a norm; where do norms come from; how do they change; how do they influence behavior; and under what conditions are they likely to become powerful? Their central theoretical contribution is the ``norm life cycle,'' a stylized three-stage account of norm emergence, norm cascade, and internalization. Different actors, motives, and mechanisms dominate at each stage. A second major contribution is the claim that norms and rationality should not be treated as opposites. Political actors often behave strategically while attempting to alter social understandings, identities, and preferences. Finnemore and Sikkink call this process ``strategic social construction.''

\subsection{What is a norm?}

A norm is defined as a standard of appropriate behavior for actors with a given identity. The concept is closely related to the sociological idea of an institution, but the distinction is one of aggregation: a norm is an individual rule or standard, whereas institutions are collections of interrelated rules and practices. This distinction matters because broad phenomena such as sovereignty or slavery are not single norms but institutional complexes composed of multiple norms.

The authors also distinguish regulative norms, which constrain and order behavior, from constitutive norms, which create new actors, interests, roles, or categories of action. Yet all norms have a prescriptive dimension: they contain a sense of ``oughtness.'' Norms are not merely regular patterns of behavior. They carry shared moral evaluation, so conformity can generate approval and violation can generate criticism, stigma, or demands for justification. When a norm is deeply internalized, conformity may produce no visible reaction because the behavior is taken for granted.

This prescriptive quality provides researchers with an empirical strategy. Norms cannot be directly observed, but they leave traces in political discourse. Actors explain and justify behavior precisely when a standard of appropriateness exists. For example, if a state feels compelled to justify continued use of a prohibited weapon, that justification itself is evidence that a norm is politically present. Crucially, researchers must distinguish evidence of a norm's existence from the behavior the norm is supposed to explain; otherwise the argument becomes circular.

Norms vary in strength and scope. Some are global, others regional; some are weakly shared, others nearly universal. The important question is therefore not simply whether a norm exists, but how agreement around it develops. This leads to the life-cycle framework.

\subsection{Domestic and international norms}

Domestic and international norms are deeply intertwined. Many international norms begin as domestic demands and are later internationalized by activists and political entrepreneurs. Women's suffrage is a central example: what began as a domestic claim in a small number of countries eventually became an international standard of legitimate political membership.

International norms also operate through domestic political structures. Their effects are filtered by local institutions, organizations, and competing norms. Domestic norm entrepreneurs can use international standards as resources in national political struggles, producing what the authors describe as a two-level norm game. This interaction is especially important in the early phases of a norm's life, before it has become strongly institutionalized internationally. Once a norm has cascaded and become widely accepted, international social pressure can become more important than domestic mobilization in generating new adopters.

\subsection{Norms as sources of stability and change}

Constructivist theory is often strong at explaining social order. Shared ideas and standards of appropriateness regularize behavior and create stable expectations. Yet this raises the same challenge that realism faces with material structure: how does structural change occur? For realists, major changes in the distribution of power can transform the system; for constructivists, shifts in norms and ideas can play an analogous role.

Finnemore and Sikkink therefore turn from the stabilizing effects of norms to the microprocesses that transform them. The life-cycle model identifies recurring mechanisms through which new standards emerge, spread, and eventually become taken for granted.

\subsection{Stage 1: norm emergence}

The first stage is driven primarily by norm entrepreneurs operating through organizational platforms. Norm entrepreneurs are actors who have strong convictions about appropriate behavior and actively attempt to persuade others to accept new standards. Norms do not simply appear; they are built through political agency.

Entrepreneurs are important because they name, interpret, and dramatize problems. This process of ``framing'' constructs new ways of understanding an issue. Successful frames resonate with broader social understandings and redefine what actors see as morally or politically appropriate. New norms always arise in a contested normative environment, not in a vacuum. Advocates must confront existing standards and interpretations. Henry Dunant and the early Red Cross movement, for example, had to persuade military authorities that wounded soldiers and medical personnel should be protected rather than treated as military resources. Suffrage activists had to challenge entrenched conceptions of women's proper social and political roles.

Because existing standards define what is currently appropriate, entrepreneurs may need to behave deliberately ``inappropriately'' to challenge them. Civil disobedience, hunger strikes, public protest, and other disruptive tactics can dramatize the inconsistency between prevailing practices and new moral claims. At this stage, a simple appeal to a logic of appropriateness cannot explain behavior because the content of appropriateness is exactly what is under dispute.

The motivations of norm entrepreneurs are often difficult to reduce to material self-interest. Finnemore and Sikkink emphasize empathy, altruism, and ideational commitment. Some advocates act because they identify with the suffering or rights of others, because they believe in the intrinsic value of a principle, or because they redefine their own interests through new moral understandings. The point is not that entrepreneurs are irrational, but that their goals cannot always be specified in purely material terms.

Organizational platforms are equally important. Entrepreneurs require organizations through which they can mobilize resources, produce information, reach audiences, and pressure states. Sometimes organizations are created specifically to promote a norm, as with the Red Cross or many advocacy NGOs. In other cases entrepreneurs operate from existing international organizations such as the UN, World Bank, or International Labour Organization. The structure, staffing, professional culture, and mandate of these organizations influence which norms they promote and how they frame them.

Expertise and professional knowledge are powerful organizational resources. Professionals help define problems and legitimate specific solutions. At the same time, professional backgrounds can block norms that do not fit dominant disciplinary assumptions. The content of organizationally promoted norms therefore reflects internal structures as well as the ideas of individual entrepreneurs.

\subsection{Institutionalization and the tipping point}

To move beyond the emergence stage, norm entrepreneurs generally need to persuade a critical mass of states to become norm leaders. Institutionalization in treaties, international organizations, or national law often helps by clarifying the norm, specifying violations, and establishing procedures for coordinating pressure against norm breakers. Institutionalization is not strictly necessary for a cascade, but it can make diffusion easier and more coherent.

The shift from emergence to cascade occurs at a tipping point. Finnemore and Sikkink cannot specify an exact universal threshold, but empirical studies suggest that tipping rarely occurs before roughly one-third of states have adopted the norm. The identity of adopters matters as much as their number. Some states are ``critical'' because the substantive objective of the norm cannot be achieved without them. Major producers or users of land mines, for example, are more normatively consequential for a land-mine ban than states that never use the weapon. Other states may be important because of prestige or moral authority.

Women's suffrage illustrates the threshold dynamic. Before around 1930, adoption was strongly associated with domestic suffrage mobilization. After roughly twenty states had adopted women's suffrage, international and transnational pressures became much more important, and adoption accelerated even in countries without strong domestic movements. Similarly, support for the land-mine ban reached around one-third of states shortly before the rapid cascade that produced the Ottawa Treaty.

\subsection{Stage 2: norm cascade}

After the tipping point, the dominant mechanism shifts from entrepreneurial persuasion to socialization. Norm leaders, international organizations, and transnational networks pressure additional states to conform. The process includes emulation, praise, criticism, diplomatic pressure, reputational incentives, and sometimes material rewards or sanctions. Calling this simply ``contagion'' is misleading because diffusion is actively produced by actors engaged in socialization.

States adopt cascading norms partly because international society shapes state identity. Once enough relevant actors accept a new standard, conformity becomes part of what it means to be a legitimate member of a particular social category, such as a modern state, liberal democracy, or European state. Three motives are especially important: international legitimation, reputation or esteem, and conformity. Governments may follow norms because they want external recognition, because leaders care about their self-image and standing, or because deviation becomes socially costly.

Socialization can therefore transform the meaning of interests. A government may initially comply for reputational reasons, but repeated conformity can gradually reshape domestic institutions and identities. At this stage, international rather than domestic pressure often becomes the main source of further diffusion.

\subsection{Stage 3: internalization}

At the final stage, norms become so widely accepted that they acquire a taken-for-granted character. Actors no longer seriously debate whether the norm should be followed. Women's suffrage, the prohibition of slavery, and the protection of medical personnel in war are examples. Internalized norms are powerful precisely because compliance no longer appears to be a choice requiring justification.

Internalization is promoted by law, bureaucratic routines, professional training, and repeated behavior. Professions are especially important because they socialize individuals into standards of appropriate conduct. Lawyers, doctors, economists, soldiers, and other professional groups carry normative commitments into state bureaucracies and international organizations. Repeated interactions can also generate habits that gradually change identities and preferences, as neofunctionalist accounts of European integration suggested.

This stage is empirically difficult to study because the strongest norms are often nearly invisible. Political debate focuses on contested rules, while taken-for-granted standards disappear into the background of routine behavior. Sociological institutionalists have been particularly useful in ``denaturalizing'' such norms and showing how apparently obvious practices such as sovereignty, markets, and individualism are historically constructed.

\subsection{Which norms matter?}

Finnemore and Sikkink identify several conditions that may make norms more influential. First is legitimation. States whose domestic elites face legitimacy problems or whose international status is insecure may be especially receptive to internationally approved norms because adoption can enhance their standing.

Second is prominence. Norms associated with states or actors perceived as legitimate, successful, or admirable may diffuse more readily. Conversely, norms associated with defeated states or failed political projects may become discredited.

Third are intrinsic characteristics of the norm. Norms protecting vulnerable groups from bodily harm and norms grounded in equality may possess broad cross-cultural resonance. This may help explain the relative success of campaigns against slavery, land mines, and violence against women. However, the authors are cautious about teleological claims that certain ``good'' norms must inevitably prevail. Political construction and historical contingency remain central.

Fourth is adjacency or fit with existing norms. New claims are often more persuasive when they can be grafted onto established standards. Legal systems make this especially explicit, since arguments gain force through analogy and precedent. New norms are more likely to spread when they can be presented as logical extensions of already accepted principles.

Finally, ``world time'' matters. Historical shocks such as wars, crises, and the collapse of political orders can discredit old norms and create openings for alternatives. Growing international organization, communication, and interdependence may also accelerate normative change. The authors note that women's suffrage required many decades to move through emergence and cascade, while later norms concerning violence against women developed much more rapidly.

\subsection{Norms, rationality, and strategic social construction}

A major theoretical argument of the article is that the opposition between norms and rationality is misleading. Rational-choice analysis does not require actors to value only material goods; utilities can include social status, identity, legitimacy, moral commitments, and other ideational objectives. At the same time, constructivists should not assume that norm-following behavior is always nonstrategic.

Norm entrepreneurs often act strategically. They select frames, audiences, organizational platforms, and institutional venues in order to change other actors' preferences and identities. States may adopt norms strategically to improve legitimacy or reputation. Actors can therefore use rational means to pursue the transformation of the social context in which future rational choices will be made. This is strategic social construction.

The relationship between a logic of consequences and a logic of appropriateness varies across the life cycle. Early entrepreneurs can calculate strategically while pursuing principled goals. During cascades, states may conform for reputation, legitimacy, or esteem. At the internalization stage, behavior becomes more habitual and less consciously chosen. The same norm can thus be associated with different motivational logics at different moments.

The authors also stress that rationality itself is socially embedded. Actors do not choose preferences, identities, and strategies in a vacuum; norms define legitimate goals, available roles, and acceptable means. Conversely, strategic actors continually reproduce and sometimes transform these normative structures. International politics is therefore characterized by the mutual constitution of agency and social structure rather than a clean separation between rational and normative action.

\subsection{Conclusion}

Finnemore and Sikkink's framework explains normative change as a patterned but nonautomatic political process. Norm entrepreneurs use persuasion and organizational platforms to create new standards; once a critical mass of important actors adopts them, socialization produces a cascade; and after widespread institutionalization, some norms become internalized and taken for granted. The likelihood and speed of success depend on legitimacy, the identity of norm leaders, the substantive content and fit of the norm, historical context, and organizational structures.

The broader theoretical lesson is that normative and rationalist explanations need not be rivals. Political actors frequently strategize within normative environments and strategically attempt to change those environments. A comprehensive explanation of international change must therefore analyze both the social construction of interests and identities and the purposive action of agents operating within, reproducing, and transforming normative structures.

\section{Lesch --- From Norm Violations to Norm Development: Deviance, International Institutions, and the Torture Prohibition}

\subsection{The puzzle and theoretical argument}

Max Lesch asks how violations affect international norms. The torture prohibition is an especially demanding test because it is formally absolute and non-derogable yet has repeatedly been violated, including by democratic states. One might expect repeated violations to weaken the norm by demonstrating that states do not comply. Lesch instead argues that violations can actively develop and sometimes strengthen norms because international institutions determine whether contested practices count as deviance and then incorporate those determinations into formal and informal lawmaking.

The central insight is that a norm violation is not simply an objective fact waiting to be observed. International actors disagree about what happened, how the relevant rule should be interpreted, and whether a particular action crosses the threshold of prohibited behavior. Courts, commissions, expert committees, and other ``norm-applying institutions'' therefore play a constitutive role. They investigate facts, interpret norms, label conduct as deviant or compliant, and publicize their conclusions. These applications subsequently become precedents or inputs into treaties, general comments, and other legal instruments. Violations thus help specify the meaning and scope of norms.

Lesch develops a two-pronged model. The first component is norm application: institutions conduct fact-finding and determine deviance by linking factual findings to normative interpretations. The second component is lawmaking: case-specific decisions feed into both formal and informal processes that reaffirm, clarify, or alter the norm. Norm development is therefore cyclical. New interpretations influence future cases, which can produce additional challenges and further reinterpretation. Importantly, development is not necessarily linear or progressive; institutional applications can clarify norms but can also generate ambiguity.

\subsection{Norm application and lawmaking}

Norm application has both empirical and normative dimensions. Institutions first establish the facts of an alleged violation. They may conduct on-site investigations, collect documents, hear witnesses, review state and NGO reports, and evaluate competing accounts. Courts often rely on adversarial procedures, while treaty bodies commonly use expert review. They then interpret the relevant legal norm and decide whether the established facts constitute conformity or deviance.

The resulting decision gives meaning to both the specific behavior and the norm itself. Calling an act illegal aggression rather than self-defense, or torture rather than lesser ill-treatment, defines the boundaries of acceptable conduct for future cases. Institutions can therefore ``fix'' contested meanings even when their authority is limited and even when states dispute the result.

Lesch distinguishes case-specific and general lawmaking and formal and informal outputs. Court judgments and advisory opinions are formal, case-specific forms. Treaty-body observations and reports are informal and case-specific. Treaties are general and formal, while general comments and standard-setting documents are general and informal. All can contribute to norm development by creating reference points that later institutions and states invoke as precedents.

\subsection{European human-rights institutions: Greece and the United Kingdom}

The first comparison concerns the European Commission on Human Rights and the European Court of Human Rights in the 1960s and 1970s. In the Greek case, Scandinavian states brought complaints against the military junta after reports of systematic torture. The Commission gathered written evidence, heard witnesses, and visited Greece. It found practices including beatings, electric shocks, falanga, mutilation, and sensory deprivation. The Commission then developed one of the first international definitions of torture in human-rights law, treating torture as deliberate, aggravated inhuman treatment that causes severe physical or mental suffering for purposes such as obtaining information or punishment. Greece was found in violation.

The second major case was Ireland v. United Kingdom, concerning British use of the ``five techniques'' in Northern Ireland: wall-standing, hooding, white noise, sleep deprivation, and deprivation of food and drink. The Commission found that the techniques intentionally caused severe suffering and met the threshold of torture. The European Court later agreed that Article 3 had been violated but reclassified the techniques as inhuman and degrading treatment rather than torture. It introduced the idea that torture carried a ``special stigma'' and required particularly serious and cruel suffering.

This judicial move created ambiguity. Rather than simply strengthening the prohibition, the two European institutions produced competing understandings of where the torture threshold lay. The Court's narrow interpretation later had permissive consequences because governments could cite it to argue that coercive interrogation techniques remained below the threshold of torture.

At the same time, these cases became important instances of lawmaking. The Commission's definition in the Greek case influenced the 1975 UN Declaration on Torture. When states later drafted the 1984 Convention Against Torture, they debated precisely the ambiguity created by the European cases. The final Convention defined torture around intentional infliction of severe physical or mental pain or suffering for specified purposes by or with the acquiescence of public officials, while omitting the especially restrictive language about ``aggravated'' ill-treatment. Thus earlier violations and the institutions' efforts to classify them directly shaped the formal international definition of torture.

\subsection{The Committee Against Torture: Israel and the United States}

The second comparison concerns the UN Committee Against Torture, a quasi-judicial expert body that reviews state compliance with the Convention Against Torture. In the 1990s and early 2000s, the Committee confronted interrogation techniques used by Israel and the United States in counterterrorism operations. Both governments advanced interpretations that sought to preserve room for coercive practices and invoked legal ambiguities concerning torture, cruel and degrading treatment, territorial scope, and exceptional security circumstances.

The Committee rejected these efforts. In its review of Israel, it challenged the concept of ``moderate physical pressure'' and treated the relevant interrogation practices as inconsistent with the Convention. In the U.S. case, it rejected arguments that certain detention sites or techniques fell outside the Convention's protections. It specifically identified practices such as waterboarding, sexual humiliation, short shackling, and the use of dogs to induce fear as prohibited forms of torture or cruel, inhuman, and degrading treatment.

These conclusions were not legally binding judgments in the same sense as a court decision, but they still mattered. They publicly labeled Israel and the United States as deviators, provided legal resources to domestic and international critics, and created an authoritative counterweight to attempts to narrow the torture prohibition. Israel's High Court later ruled that moderate physical pressure was unlawful, while the Obama administration rescinded the U.S. ``enhanced interrogation'' program. The United States subsequently returned to the Committee seeking to distance itself from the deviator label and acknowledged that some of the practices constituted torture.

\subsection{General Comment No. 2 and norm reinforcement}

The Israeli and U.S. disputes also stimulated general informal lawmaking. In 2008 the Committee adopted General Comment No. 2 on Article 2 of the Convention. The document clarified that the obligations to prevent torture and other cruel, inhuman, or degrading treatment overlap substantially and that neither can be weakened through exceptional security arguments. It reaffirmed the absolute and non-derogable nature of the prohibition and specified that it applies during terrorism, violent crime, and international or noninternational armed conflict.

General Comment No. 2 also clarified territorial scope, stating that the Convention applies wherever a state exercises de jure or de facto control. This directly countered attempts to exclude secret detention facilities, occupied territories, or extraterritorial counterterrorism operations. In drafting the comment, Committee members explicitly drew on the disputes with Israel and the United States. Thus concrete norm applications were generalized into a broader interpretation of the treaty.

Although the United States contested the legal authority of the General Comment, the document gained wider recognition. The UN General Assembly referenced it, and human-rights courts subsequently cited it. Lesch therefore treats it as evidence that informal institutional lawmaking can become part of the wider development of international law even without the formal authority of treaty amendment or binding adjudication.

\subsection{Conclusion and contribution}

The comparison shows that norm violations can be productive moments in normative development. Violations expose ambiguities, force institutions to interpret rules, and generate records of deviance that can later shape formal and informal lawmaking. The Greek and British cases helped construct the modern international definition of torture, though the European Court also introduced ambiguity. The Israeli and U.S. cases prompted the Committee Against Torture to reaffirm the norm's absolute and non-derogable character and to clarify its territorial and substantive scope.

Lesch's argument therefore moves beyond approaches that treat international institutions merely as arenas in which states contest norms. Institutions themselves exercise agency by determining deviance and fixing meanings. Their influence depends on fact-finding practices, professional interpretation, precedents, state recognition, and the interaction between specific decisions and general lawmaking. Persistent violations do not automatically erode a norm. Under some conditions, the institutional response to violations can clarify and reinforce it. At the same time, because interpretations remain contested, norm development is neither automatic nor necessarily progressive.

\section{Mitchell and Carpenter --- Norms for the Earth: Changing the Climate on ``Climate Change''}

\subsection{The problem: climate governance dominated by interests}

Ronald B. Mitchell and Charli Carpenter argue that climate change is a ``super-wicked'' collective-action problem characterized by severe asymmetries. The states most vulnerable to climate impacts often have the least capacity to reduce global emissions, while some of the most powerful and carbon-intensive states face weaker immediate incentives to accept costly mitigation. This structure gives major emitters strong veto power and makes ambitious agreements difficult.

Most climate governance has addressed this problem through what March and Olsen call a logic of consequences. Governments and institutions try to persuade states that mitigation serves their material interests by providing scientific information, creating reciprocal commitments, offering finance and technology transfers, developing emissions trading, monitoring compliance, or altering the economic costs of carbon. The underlying logic is instrumental: states should reduce emissions because the benefits exceed the costs or because others will reciprocate.

Mitchell and Carpenter do not reject these tools, but they argue that the record of climate governance shows their limits. Global emissions continued to rise despite decades of scientific assessment, negotiation, and increasingly sophisticated policy instruments. Strategies that are politically acceptable to powerful emitters are often too weak to address the problem, while strategies large enough to prevent severe climate impacts often appear materially costly and therefore politically unacceptable.

The authors propose that climate advocates should draw lessons from human-security and humanitarian-disarmament campaigns, where norm entrepreneurs have sometimes overcome similarly unfavorable distributions of interests by shifting political debate from material calculation to ethics. A logic of appropriateness asks not ``What maximizes my material payoff?'' but ``What is the right thing for an actor like me to do?'' States may accept costly behavior because it is tied to a valued identity, moral obligation, or conception of civilized conduct.

\subsection{Interest-based and ethics-based governance}

Interest-based and ethics-based logics are ideal types rather than mutually exclusive categories. Trade regimes rely heavily on reciprocity and exchange, while human-rights regimes more explicitly ask states to respect principles even when compliance is costly. Arms control and environmental governance often combine the two. Chemical and biological weapons regimes, for example, involve monitoring and strategic incentives but also strong moral stigmas. The whaling regime evolved from a resource-management framework toward an ethical discourse that increasingly treated whaling itself as morally problematic.

Climate governance, by contrast, remains dominated by scientific and economic reasoning. The UN Framework Convention and the Paris Agreement recognize vulnerable groups, sustainability, poverty, and ecosystems, but tend to frame these concerns as contexts or constraints on climate policy rather than as the primary moral reason for deep emissions reductions. Ethical climate arguments exist, as illustrated by Pope Francis's encyclical on care for the common home, but they have not become the dominant logic of international negotiations.

The authors argue that interest-based persuasion is especially weak when large emitters simply do not believe deep reductions serve their material interests. Under these circumstances, changing prices and information may be insufficient. Political strategy must also change the social meaning of high-carbon behavior and the identity costs of refusing action.

\subsection{Five ethics-based strategies}

Drawing on research on human rights and humanitarian disarmament, Mitchell and Carpenter identify five strategies that have helped norm entrepreneurs overcome resistant veto players.

First, advocates can promote discursive shifts. They can reframe a policy problem from costs and benefits to rights and wrongs. Successful humanitarian campaigns changed the terms of debate by linking disputed practices to accepted metanorms such as bodily integrity, civilian protection, and the prohibition of indiscriminate harm. Norm ``grafting'' connects a new claim to older, already legitimate principles. Activists can also reverse the burden of proof, making continued participation in a harmful practice something that requires public justification. In the nuclear-weapons-ban campaign, advocates reframed nuclear weapons from instruments of national security and deterrence into a humanitarian threat whose use would impose catastrophic and indiscriminate suffering.

Second, advocates can define and promote norm-related identities. Norms influence behavior when they become linked to identities states value, such as being a responsible international citizen, humanitarian leader, or law-abiding democracy. Norm entrepreneurs can praise states that perform these roles and frame resistance as inconsistent with the identity a state claims. In humanitarian disarmament, middle powers have often sought status as norm leaders, and campaigns have used this desire to encourage stronger commitments.

Third, activists can mobilize pride and shame. Naming and shaming works by defining a behavior as morally inappropriate and compelling governments to explain themselves in the language of the norm. Praising early adopters can be just as important, especially in the emergence phase. Social pressure can influence even nonparties to treaties. The land-mine and cluster-munitions regimes affected production and use by states outside the treaties because the weapons acquired a stigma that altered reputational incentives.

Fourth, advocates can use ethical frames to mobilize transnational networks. The success of norm campaigns depends not merely on the number of organizations involved but on network structure. Influential gatekeepers, brokers, and central coordinating organizations can help create a coherent message and connect previously separate constituencies. The International Campaign to Abolish Nuclear Weapons benefited from the participation of the International Committee of the Red Cross and from organizations that linked nuclear disarmament to the established humanitarian-disarmament community. This helped move the issue from a relatively isolated peace movement into a broader and more professional humanitarian network.

Fifth, norm entrepreneurs can shift institutional forums to marginalize veto players. Consensus rules often allow powerful opponents to block stronger norms. Campaigns against land mines, cluster munitions, and nuclear weapons succeeded partly by moving negotiations to venues in which committed norm leaders could proceed without the consent of major military powers. Venue shifting changes the political agenda, rewards leaders, exposes laggards, and allows a strong ethical principle to be codified even if important states refuse to join. The resulting treaty can still reshape discourse and behavior beyond its formal membership.

\subsection{The nuclear weapons ban as a model}

The 2017 Treaty on the Prohibition of Nuclear Weapons is the article's most important illustration. Earlier nuclear regimes had emphasized deterrence, arms control, and nonproliferation while implicitly accepting the continued possession of nuclear weapons by a small group of states. Ban advocates deliberately challenged this framework by emphasizing catastrophic humanitarian consequences and the moral illegitimacy of nuclear use and possession.

A series of humanitarian-impact conferences privileged this discourse and reduced the agenda-setting power of nuclear-armed veto players. The process did not require consensus, so opponents could not block negotiations. Advocacy networks increasingly presented nuclear possession as inconsistent with responsible humanitarian identity. The treaty was eventually adopted by 122 states despite explicit rejection by major nuclear powers.

For Mitchell and Carpenter, the importance of this outcome is not limited to legal membership. Humanitarian bans can create stigma and change the rhetorical balance of power even when powerful states remain outside them. Once a prohibitionary norm exists, nonparticipants must explain why their behavior remains legitimate. Similar dynamics have affected land mines and cluster munitions. This suggests that an ethically framed climate norm could influence major emitters even before universal participation is achieved.

\subsection{Applying these strategies to climate politics}

The first recommendation is to delegitimize the assumption that climate action must always be justified as economically beneficial to the emitter. When governments claim that ambitious mitigation is too costly, advocates often respond with lower cost estimates or new material benefits. This preserves the logic that action is required only when it pays. An ethics-based approach would instead argue that states have obligations to avoid imposing severe harm on vulnerable populations, future generations, other species, and threatened territories even when doing so requires sacrifice.

Norm grafting could connect climate responsibility to established rules protecting human rights, territorial integrity, indigenous rights, and bodily security. Sea-level rise threatens the territorial existence of states, while climate-related mortality and displacement implicate widely accepted rights. Such linkages could redefine excessive emissions as conduct requiring moral and political justification.

Climate policy can also activate identities. Kyoto and Paris already create standards against which states describe themselves as responsible or irresponsible actors. The Paris Agreement's expectation that nationally determined contributions become progressively more ambitious provides an example of rhetorical entrapment: once governments voluntarily accept the identity of states making ``ambitious and progressive'' commitments, future backsliding becomes harder to justify. The voluntary nature of these pledges may paradoxically strengthen their normative significance because states cannot easily claim that obligations were externally imposed.

Pride and shame could be used more systematically. States, cities, firms, and other actors that adopt ambitious policies can be publicly recognized as norm leaders, while persistent high emitters can be framed as laggards. The goal would not merely be to calculate the price of carbon but to change the social meaning of carbon-intensive behavior.

Transnational climate networks might also learn from humanitarian campaigns by strengthening coordination around a clear ethical frame. Climate advocacy is broad and fragmented, encompassing environmentalism, development, labor, indigenous rights, religion, science, and economic policy. A successful normative campaign would need influential gatekeepers and brokers capable of connecting these constituencies around a shared claim about the moral unacceptability of imposing severe climate harm on vulnerable others.

Finally, venue shifting may help when consensus-based climate institutions give veto players excessive control. Smaller coalitions of willing states could develop strong normative standards in alternative forums and then use those standards to generate stigma and pressure outside the coalition. The objective would not necessarily be to replace the UN climate regime, but to create supplementary arenas in which more ambitious ethical norms can emerge without being diluted by actors committed to weaker action.

\subsection{Conclusion}

Mitchell and Carpenter argue that climate politics should not abandon material incentives, scientific expertise, or institutional bargaining, but should add a more ambitious ethics-based strategy. The central weakness of the prevailing approach is its assumption that major emitters must first be persuaded that deep mitigation is in their material self-interest. Human-security campaigns show that political behavior can instead be changed by redefining what responsible actors are expected to do.

The five strategies---discursive reframing, identity construction, pride and shame, transnational network mobilization, and forum shifting---offer practical ways to build a logic of appropriateness around climate action. Their success in humanitarian disarmament does not prove that they will automatically work for climate change, which differs in important ways from weapons bans. But they provide alternative models worth testing. The authors therefore call for both practitioners and scholars to devote greater attention to the deliberate construction of strong ethical norms that treat dangerous emissions not simply as economically inefficient, but as socially inappropriate and morally unacceptable behavior.
